\section{Related Work}
\label{sec:relwork}
The gap this chapter documents can be stated before the evidence: graph reduction is a
developed literature with surveys of its own
\citep{hashemi2024survey,shabani2023summarization}, circuit representation learning is a
developed literature with a growing model family, and no published work applies the first
to whole AIGs as input reduction for the second. The subsections below establish each
half, then the two nearest attempts at a bridge, and~\ref{sec:relwork:gap} returns to the
claim once the evidence for it is on the table.

\subsection{Machine Learning for Logic Synthesis}
\label{sec:relwork:ml4eda}

\subsubsection{Quality-of-Result and Optimizability Prediction}
\label{sec:relwork:ml4eda:qor}
The direct precedents predict synthesis Quality of Result (QoR) from a circuit before
running the tool. OpenABC-D \citep{chowdhury2021openabc} is the reference dataset for the
task and publishes, among its graph-level labels, the percentage of nodes a recipe
optimizes away, essentially the target~\eqref{eq:prelim:target}; that overlap is what
makes its accompanying SynthNet a fair baseline here rather than a repurposed model
(\ref{sec:relwork:ml4eda:baselines}). Recipe-conditioned models jointly encode the
circuit and the synthesis sequence \citep{qor_transformer2022,wu2022lostin}. One
terminological collision is worth heading off: LOSTIN \citep{wu2022lostin} attaches a
super-node that encodes the synthesis \emph{sequence}, a temporal use of the same device
this thesis uses structurally, and the two share nothing beyond the name.

\subsubsection{Evaluation Protocols and the Unseen-Design Gap}
\label{sec:relwork:ml4eda:protocols}
``Unseen'' means three different things in this literature: an unseen recipe, an unseen
design, or an unseen design-recipe pair whose parts were both seen, and the third is far
easier than the second. Generalization to unseen designs is the stated motivation of much
of the field, yet matched seen and unseen results are rare: of the papers surveyed for
this thesis, 4 of 63 report both sides, among them OpenABC-D \citep{chowdhury2021openabc}
and LOSTIN \citep{wu2022lostin}, while HOGA \citep{deng2024hoga} motivates on the gap
without publishing a seen-design baseline, so its degradation cannot be computed.
% TODO(citation): LSOformer and Jiang et al. are the other two; add keys before naming them.
The canonical unseen-IP split trains on 16 small designs and tests on 8 large ones,
confounding design novelty with size extrapolation. Reported error moves more with the
benchmark than with the architecture, from a few percent on small homogeneous suites to
above twenty percent on OpenABC-D, and within one test set the error concentrates in a
few designs rather than spreading uniformly. Those three observations motivate RQ1a's
matched three-protocol comparison and the per-design reporting used throughout
(\ref{sec:intro:rq1}).

\subsubsection{Learned Synthesis Script Generation}
\label{sec:relwork:ml4eda:scripts}
A separate line learns to construct or select synthesis scripts, by reinforcement
learning or sequence prediction over the same primitive transformations. It is the
downstream application the motivation invokes, and it is out of scope here
(\ref{sec:intro:scope:prediction}): script generation consumes exactly the kind of cheap
optimizability estimate this thesis builds, which is why the prediction task is worth
solving on its own.
% TODO(citation): add one representative script-generation reference.

\subsubsection{Circuit Representation Learning}
\label{sec:relwork:ml4eda:representation}
The DeepGate line \citep{li2022deepgate,shi2024deepgate3,deepgate4_2025} learns
general-purpose gate embeddings on AIGs, and PolarGate \citep{polargate2024} identifies
the handling of inverted connections as the representational bottleneck of AIG encoders,
external support for treating inversion as a relation summarization must preserve rather
than average away (\ref{sec:method:summary:wl}). Two facts from this line shape the
present study. First, DeepGate3 and DeepGate4 answer scale by growing the \emph{model},
with sparse attention and larger training corpora, where this thesis shrinks the
\emph{input}. DeepGate4 doubles as a baseline (\ref{sec:relwork:ml4eda:baselines}), so
the contrast is measured rather than rhetorical. Second, these models train on extracted
subcircuits of tens to a few thousand gates. Extraction is a reduction, applied at corpus
construction and never named or evaluated as one, and that unexamined step is half of the
gap this chapter closes in~\ref{sec:relwork:gap}.

\subsubsection{Models Used as Baselines}
\label{sec:relwork:ml4eda:baselines}
Three published models serve as tier-2 baselines,
and~\ref{sec:method:baselines:published} states only how each was adapted, so what each
one \emph{is} belongs here. SynthNet
\citep{chowdhury2021openabc} is a recipe-conditioned graph convolutional QoR regressor,
the closest published analogue of the label predicted here. HOGA \citep{deng2024hoga}
applies hop-wise attention to multi-hop features precomputed per node, so its depth of
field is fixed before training and its trunk never sees connectivity, a design aimed at
scalable, distributable training. DeepGate4 \citep{deepgate4_2025} runs sparse attention
over a distance-bounded virtual edge set and is the architecture-scaling counterpart to
the input scaling studied here. Each comparison licenses a different conclusion: against
SynthNet the claim is task accuracy, against HOGA it is whether connectivity at training
time matters, and against DeepGate4 it is input against architecture scaling at matched
memory.

\subsubsection{Benchmarks and Datasets for ML in EDA}
\label{sec:relwork:ml4eda:datasets}
OpenABC-D \citep{chowdhury2021openabc} remains the largest published synthesis-prediction
corpus: 29 designs, 1500 recipes each, labels per design-recipe pair. The dataset built
for this thesis (\ref{sec:method:data}) differs in axis rather than only in size: it is
tiered by optimization state, pairing each base graph with already-optimized variants of
itself (\ref{sec:method:data:tiers}), because the reduction questions RQ3 and RQ5 need
graphs at several points along the optimization trajectory, not several recipes on one
starting point.
% TODO(citation): OpenLS-DGF and the ML4EDA benchmark position paper; add or drop.
